% ===========================================================================
%  Worksheet --- Field Extensions and the Ring F[x]
%
%  THE WRITE-UP, and it is the assistant's to type.
%    The assistant transcribes each question or topic faithfully and emits an
%    empty, marked solution region beneath it. The student does the
%    mathematics; the assistant typesets what they agreed was right.
% ===========================================================================
\documentclass[11pt]{article}
\usepackage{coursemacros}

\title{Worksheet --- Field Extensions and the Ring $F[x]$}
\author{Mikey Ferguson}
\date{\today}

\begin{document}
\maketitle

% One problem/solution pair per thing worked, in the order it was assigned or
% chosen. A region stays empty until the answer in it is agreed correct.

\begin{problem}{6}
  Let $K = \mathbb{Q}(\sqrt{2},\sqrt{3})$ and $\gamma = \sqrt{2} + \sqrt{3}$.

  \begin{enumerate}
    \item[(a)] Show that $\degree{K}{\mathbb{Q}} = 4$.
    \item[(b)] Show that $K = \mathbb{Q}(\gamma)$, and compute
      $m_{\gamma,\mathbb{Q}}$.
    \item[(c)] Find three proper intermediate fields.
  \end{enumerate}
\end{problem}

% ===== SOLUTION 6(a) =====
% Transcribed from the student's slate, 2026-09-18 (page t0054-r42, turns
% leading to cards 0047-0050). Typeset exactly as written: same steps, same
% order, nothing added.
% Two reason-corrections, both agreed on cards 0048 and 0049 and made here
% rather than left to the student:
%   - the page justifies ab = 0 with "since a,b in Q"; the reason is that
%     sqrt(2) is irrational.
%   - the page flags each branch with "since b in Q" / "since a in Q"; the
%     contradictions are the irrationality of sqrt(6) and of sqrt(3).
\begin{solution}
  \textbf{(a)} First, $\sqrt{3} \notin \mathbb{Q}(\sqrt{2})$.

  Suppose $\sqrt{3} = a + b\sqrt{2}$ with $a, b \in \mathbb{Q}$. Squaring,
  \[
    3 = a^{2} + 2ab\sqrt{2} + 2b^{2},
  \]
  so $2ab\sqrt{2} = 3 - a^{2} - 2b^{2}$ is rational, and since $\sqrt{2}$ is
  irrational this forces $ab = 0$.

  \emph{Suppose $a = 0$.} Then $\sqrt{3} = b\sqrt{2}$, so $2b = \sqrt{6}$. With
  $b \in \mathbb{Q}$ that makes $\sqrt{6}$ rational, and it is not.

  \emph{Suppose $b = 0$.} Then $a = \sqrt{3}$. With $a \in \mathbb{Q}$ that
  makes $\sqrt{3}$ rational, and it is not.

  So $\sqrt{3} \neq a + b\sqrt{2}$ for any $a, b \in \mathbb{Q}$, that is,
  $\sqrt{3} \notin \mathbb{Q}(\sqrt{2})$.

  \medskip
  \textbf{Note.} $\degree{\mathbb{Q}(\sqrt{2})}{\mathbb{Q}} = 2$, since
  $\mathbb{Q}(\sqrt{2}) = \mathbb{Q}[\sqrt{2}]$ has basis $\set{1, \sqrt{2}}$
  over $\mathbb{Q}$.

  \medskip
  \textbf{Further.} $\degree{K}{\mathbb{Q}(\sqrt{2})} = 2$, since
  $K = \mathbb{Q}(\sqrt{2})[\sqrt{3}]$ has basis $\set{1, \sqrt{3}}$ over
  $\mathbb{Q}(\sqrt{2})$: the minimal polynomial of $\sqrt{3}$ there divides
  $x^{2} - 3$, so has degree $2$ or $1$, and degree $1$ is ruled out by
  $\sqrt{3} \notin \mathbb{Q}(\sqrt{2})$.

  \medskip
  So by the tower law
  \[
    \degree{K}{\mathbb{Q}}
      = \degree{K}{\mathbb{Q}(\sqrt{2})}\,
        \degree{\mathbb{Q}(\sqrt{2})}{\mathbb{Q}}
      = 2 \cdot 2 = 4 .
  \]
\end{solution}

% ===== SOLUTION 6(b) =====
% Transcribed from the student's slate, 2026-09-18 (pages t0054-r43, t0054-r44,
% t0054-r45 and t0054-r46, cards 0051-0054). Typeset exactly as written: same
% steps, same order, nothing added. Both halves are now agreed.
% One reason-correction, agreed on cards 0051 and 0053 and made here rather than
% left to the student: r43 justifies minimality with "degree 4, monic, so this
% is the minimal polynomial". Monic of degree 4 is not enough on its own. The
% reason is the degree count the student then gave on r46, and that is what the
% write-up says.
\begin{solution}
  \textbf{(b)} $(\sqrt{2}+\sqrt{3})(\sqrt{3}-\sqrt{2}) = 3 - 2 = 1$, so
  $\sqrt{3}-\sqrt{2} = \gamma^{-1}$, and hence
  $\sqrt{3}-\sqrt{2} \in \mathbb{Q}(\gamma)$.

  Take $\gamma + \gamma^{-1} = 2\sqrt{3} \in \mathbb{Q}(\gamma)$. Since
  $2 \in \mathbb{Q}$, that puts $\sqrt{3} \in \mathbb{Q}(\gamma)$. So
  $\gamma - \sqrt{3} = \sqrt{2} \in \mathbb{Q}(\gamma)$.

  So $\mathbb{Q}(\gamma) = \mathbb{Q}(\sqrt{2},\sqrt{3})$. Both contain
  $\sqrt{2}+\sqrt{3}$, $\sqrt{2}$ and $\sqrt{3}$; $\mathbb{Q}(\gamma)$ is the
  smallest field extension of $\mathbb{Q}$ containing $\sqrt{2}+\sqrt{3}$, and
  $\mathbb{Q}(\sqrt{2},\sqrt{3})$ is the smallest field extension of
  $\mathbb{Q}$ containing $\sqrt{2}$ and $\sqrt{3}$. So they must be the same.

  \medskip
  \textbf{The minimal polynomial.} Squaring twice,
  \[
    \gamma^{2} = 5 + 2\sqrt{6}
    \implies (\gamma^{2}-5)^{2} = 24
    \implies \gamma^{4} - 10\gamma^{2} + 25 = 24
    \implies \gamma^{4} - 10\gamma^{2} + 1 = 0 .
  \]

  Now $\degree{\mathbb{Q}(\gamma)}{\mathbb{Q}} = 4$, since
  $\mathbb{Q}(\gamma) = \mathbb{Q}(\sqrt{2},\sqrt{3})$ and that degree is $4$ by
  (a). So $\deg m_{\gamma,\mathbb{Q}} = 4$, meaning $x^{4} - 10x^{2} + 1$ is the
  monic polynomial which is $m_{\gamma,\mathbb{Q}}$: any lower degree cannot have
  $\gamma$ as a root.
\end{solution}

% ===== SOLUTION 6(c) =====
% Transcribed from the student's slate, 2026-09-18 (page t0054-r47, card 0055).
% Typeset exactly as written: the three fields and, beside each, the elements of
% K it omits. Nothing added beyond naming what the page marks with arrows.
\begin{solution}
  \textbf{(c)} $\degree{\mathbb{Q}(\gamma)}{\mathbb{Q}} = 4$, since
  $\mathbb{Q}(\gamma) = \mathbb{Q}(\sqrt{2},\sqrt{3})$.

  Intermediate fields of $\mathbb{Q}(\sqrt{2},\sqrt{3}) : \mathbb{Q}$:
  \[
    \begin{array}{lcl}
      \mathbb{Q}(\sqrt{2}) & \longrightarrow & \text{no } \sqrt{3},\ \sqrt{6}\\
      \mathbb{Q}(\sqrt{3}) & \longrightarrow & \text{no } \sqrt{2},\ \sqrt{6}\\
      \mathbb{Q}(\sqrt{6}) & \longrightarrow & \text{no } \sqrt{2},\ \sqrt{3}
    \end{array}
  \]

  Each is a proper intermediate field: each contains $\mathbb{Q}$ and sits
  inside $K$, each is bigger than $\mathbb{Q}$ because its generator is
  irrational, and none is all of $K$ because each omits the two elements of $K$
  named beside it. Those omissions also make the three fields distinct from one
  another.
\end{solution}

\end{document}
